% !TeX root = ../../Book5.tex
% 纲要标题：### 32.4 余切复形的统一解释
% 纲要位置：几何学计划纲要.md，第 6108 行。
% 正文状态：未撰写；本文件目前是框架。

% 纲要内容开始
% 
% * Ext 的自同构、切向及障碍次数
% * 局部到整体谱序列
% * 光滑对象的 H¹(T) 与 H²(T)
% * 半通用形式形变及 Schlessinger 条件的所选范围
% 
% 纲要内容结束

\section{余切复形的统一解释}
\label{g5:sec:3204}

待撰写

\subsection{Ext 的自同构、切向及障碍次数}
\label{g5:subsec:3204:01}
待撰写

\subsection{局部到整体谱序列}
\label{g5:subsec:3204:02}
待撰写

\subsection{\texorpdfstring{光滑对象的 \(H^1(\mathrm{T})\) 与 \(H^2(\mathrm{T})\)}{光滑对象的 H¹(T) 与 H²(T)}}
\label{g5:subsec:3204:03}
待撰写

\subsection{半通用形式形变及 Schlessinger 条件的所选范围}
\label{g5:subsec:3204:04}
待撰写
% 纲要条目排为编号小节标题，并进入目录与书签；说明段仍保留于纲要注释中。
